\documentclass[11pt,a4paper]{article}
\usepackage[utf8]{inputenc}
\usepackage[T1]{fontenc}
\usepackage[margin=1in]{geometry}
\usepackage{xcolor}
\usepackage{hyperref}

\definecolor{primarycolor}{RGB}{33,37,41}
\definecolor{secondarycolor}{RGB}{0,102,204}
\hypersetup{colorlinks=true,linkcolor=secondarycolor,urlcolor=secondarycolor}
\pagestyle{empty}

\begin{document}

Dear Hiring Manager,

I am applying for the Machine Learning Engineer -- Customer Support AI role at Quillstone Payments. My data-science experience combines machine-learning workflows, experimentation and practical engineering improvements, and I am keen to grow further into the production engineering side of machine learning.

At Thistlebox, I built a six-week recipe-demand forecasting workflow that improved WAPE from 23\% to 14\%, with an estimated 9\% reduction in food waste. I also developed weekly BigQuery churn scoring used to target retention offers; a 10\% holdout validation found churn was down 2.1 percentage points in the treated group. These projects have given me hands-on experience with Python, feature engineering, model evaluation and workflows that inform business operations.

My work so far has focused on scheduled notebook and batch-scoring workflows rather than independently deployed model APIs. I wrote the Dockerfile for a Streamlit sequential-testing app, which the data platform team hosts on Cloud Run, and since 2025 I have been writing pytest unit tests for forecasting feature code, reaching approximately 60\% coverage. I am interested in Quillstone’s focus on taking AI ideas from prototype to production in collaboration with backend engineers, data scientists and product managers, particularly to help merchants and support teams get answers with less friction.

Thank you for considering my application. I would welcome the opportunity to discuss how my experience in applied machine learning and experimentation could contribute to Quillstone’s Customer Operations AI team.

Kind regards,\\
Tobias Mwangi-Hartley

\end{document}